% Source: https://artofproblemsolving.com/community/c3h3131590_2023_emcc_guts_round_p13__p24__exeter_math_club_competition
% Stable identity: aops-topic-3131590-q24
For a prime $p\in\{2,5\}$ dividing $n$, let $v_p(n)$ be its exponent in $n$. If each such prime satisfies $nv_p(n)\ge k$, the suffix occurs infinitely often: replace $n$ by $N=n+t\,10^k\varphi(10^k)$. For a prime not dividing $n$, the base and exponent congruences preserve the power modulo $p^k$; for a prime dividing $n$, both powers are zero modulo $p^k$. Thus a finite suffix must have a solution with $nv_p(n)<k$ for at least one of these primes, implying $n<k$. If $k\le10$, the only relevant $n<k$ divisible by two or five are $2,4,5,6,8$. For $n=2,4,5,6$, unsaturation requires respectively $k\ge3,9,6,7$, but $n^n=4,256,3125,46656$ then has fewer than $k$ digits and cannot produce a suffix with nonzero leading digit. For $n=8$, its two-adic exponent is $24$, already saturated. Hence no $k\le10$ works. For $k=11$, take the suffix $10000000000=10^{10}$, attained by $n=10$. Any other solution $N$ has $Nv_2(N)=Nv_5(N)=10$, so $10\mid N$ and $N\le10$, forcing $N=10$. It is therefore blissful, and the minimum is $\boxed{11}$.
